\documentclass[10pt]{article}
\usepackage[margin=0.78in]{geometry}
\usepackage{amsmath,amssymb,booktabs}
\usepackage[T1]{fontenc}
\usepackage{lmodern}
\usepackage{microtype}
\usepackage{enumitem}
\setlist{nosep,leftmargin=*}
\title{TPC-382: A Finite \(c=1\) Origin-Family Magnitude Audit}
\author{Liang Wang\\School of Mathematics and Statistics, HUST\\Wuhan, China}
\date{September 4, 2026}
\begin{document}
\maketitle

\begin{abstract}
We audit the magnitude, rather than only the threshold profile, of a finite
prime-shell operator family used in the twin-prime exploratory line.  Two
sealed, protocol-matched \(N=2048\) panels (TPC-380 and TPC-381) provide six
origins, three \(Q\)-anchors, and four predeclared laws.  With the rule
\(\Delta=(\max-\min)/\operatorname{mean}\leq 0.01\) fixed before reading the
parent rows, all-plus is stable at all three \(Q\)'s; the full law census has
8 of 12 cells within the cap.  A separately labelled TPC-379 \(N=1024\)
scale control refutes one-percent matched-count invariance for all-plus at
high \(Q\), with a relative contrast of \(0.020813995160269608\).  These are
finite certificate results.  They do not establish source validity, an
asymptotic uniformity theorem, arithmetic power saving, or a twin-prime result.
\end{abstract}

\section{Question and scope}

The preceding origin-family replay preserved the finite band failure profile
but left magnitude stability as an explicit open diagnostic.  We ask two
finite questions:
\begin{enumerate}
\item Is the normalized band spectral magnitude stable across the two
protocol-matched \(N=2048\) origin families?
\item Is that stability common to all four declared laws, and does it persist
when the count is changed to the earlier \(N=1024\) panel?
\end{enumerate}
The first question is answered using already sealed certificates.  The second
is split into a law census and a clearly non-matched scale control; no claim is
made about an unobserved source law.

\section{Locked protocol}

For a parent panel \(P\), law \(\ell\), anchor \(Q\), and origin \(o\), let
\(s(P,\ell,Q,o)\) denote the recorded spectral value of the normalized
\(c=1\) band.  For a locked set of values define
\[
 \Delta=\frac{\max s-\min s}{|O|^{-1}\sum_{o\in O}s(P,\ell,Q,o)}.
\]
The one-percent cap is fixed at \(0.01\).  TPC-380 and TPC-381 have the same
count, block rule, shell anchors, kernel, weights, and four laws; their three
origins are combined, giving \(6\times3\times4=72\) values in 12 cells.
TPC-379 is retained only as a count-\(1024\) scale control with 36 values.

Before aggregation, the producer locks the source and certificate hashes:
TPC-379, TPC-380, and TPC-381 are respectively identified by their canonical
certificate hashes
\texttt{a41800cb}, \texttt{c80dbfab}, and \texttt{c217a475} (prefixes shown
only for readability).  The independent checker repeats the parent loading,
row-key census, and arithmetic without importing the producer.

\section{Results}

The all-plus same-count cells have relative spreads
\[
 1.9035250282068572\!\times\!10^{-5},\quad
 2.380537285421679\!\times\!10^{-5},\quad
 8.0645464844910632\!\times\!10^{-6}
\]
at \(Q=512,2048,8192\), respectively.  The high-\(Q\) mean is
\(0.66694363456350925\).  The complete same-count census is shown below;
``yes'' means that the cell's six-origin spread is at most one percent.

\begin{center}
\begin{tabular}{lccc}
\toprule
law & \(Q=512\) & \(Q=2048\) & \(Q=8192\)\\
\midrule
all-plus & yes & yes & yes\\
alternating-index & no & no & no\\
mod-4 character & no & yes & yes\\
half-split & yes & yes & yes\\
\bottomrule
\end{tabular}
\end{center}

Thus 8/12 cells pass the fixed cap.  The four failures are exactly the three
alternating-index cells and the mod-4 cell at \(Q=512\).  This is a law-
dependent magnitude census, not evidence that a law may be selected from the
data.

For the scale control, compare the mean of the six \(N=2048\) values with the
mean of the three TPC-379 \(N=1024\) values at matched law and \(Q\).  At
all-plus and \(Q=8192\),
\[
 \frac{m_{2048}-m_{1024}}{m_{1024}}
 =0.020813995160269608,
\]
which exceeds the predeclared one-percent scale-invariance cap.  Because the
two counts use different normalizations and different origin families, this
is a finite refutation of that narrowly stated hypothesis, not a theorem of
scale non-uniformity.

\section{Verification and claim firewall}

The canonical producer replay, independent aggregation replay, and 25-field
adversarial mutation suite each pass in normal and optimized Python modes.
The finite proof package therefore certifies the arithmetic transformation of
the locked parent rows and the stated finite comparisons.  The official
Session evaluator files are not present in this checkout; the repository's
Bridge-B is a local fail-closed check.

\begin{center}
\small
\begin{tabular}{p{0.43\linewidth}p{0.45\linewidth}}
\toprule
item & status\\
\midrule
parent locks and row census & PROVED\_EXACT\_FINITE\\
all-plus high-\(Q\) one-percent stability & NUMERICALLY CERTIFIED, FINITE SCOPED\\
law-dependent spread census & NUMERICALLY CERTIFIED, FINITE SCOPED\\
matched-count one-percent scale hypothesis & REFUTED, FINITE SCOPED\\
source-valid normalization, growing bound, arithmetic \(L^2\) & OPEN\\
Route-A / Route-B gates & OPEN\\
twin-prime conclusion & NONE\\
\bottomrule
\end{tabular}
\end{center}

In particular, \texttt{ARITHMETIC\_ADVANCE}=\texttt{NO} and
\texttt{FIXED\_POWER\_CREDIT}=0.  The next finite test is the predeclared
pooled cross-origin normalization audit: determine whether the observed
magnitude persistence survives a common normalization.

\end{document}
